\topic{把一条指令拆成每个时钟周期}

前文已经列出取指、译码、执行、访存和写回。现在让这些词落到一台可以运行的
机器上。教学 CPU 使用 16-bit 地址、4 个 16-bit 通用寄存器、一位零标志和
一个向低地址增长的栈。ROM 占用
\(\texttt{0x0000}\)--\(\texttt{0x7FFF}\)，RAM 占用
\(\texttt{0x8000}\)--\(\texttt{0xDFFF}\)，
\(\texttt{0xE000}\) 起始的窗口故意不连接任何设备。

\begin{pdfdiagram}
  \node[os state] (pc) {PC\\下一次取指地址};
  \node[os block,right=18mm of pc] (bus) {地址/数据/读写/ready};
  \node[os hardware,right=18mm of bus,yshift=16mm] (rom) {ROM\\指令字节};
  \node[os hardware,right=18mm of bus] (ram) {RAM\\数据与栈};
  \node[os hardware,right=18mm of bus,yshift=-16mm] (hole) {未映射窗口\\不返回响应};
  \node[os state,below=18mm of bus] (ir) {操作码与操作数};
  \node[os block,left=18mm of ir] (decode) {译码与控制};
  \node[os state,left=18mm of decode] (registers) {R0--R3、Z、SP};
  \draw[os bus] (pc) -- (bus);
  \draw[os arrow] (bus) -- (rom);
  \draw[os arrow] (bus) -- (ram);
  \draw[os arrow] (bus) -- (hole);
  \draw[os arrow] (bus) -- (ir);
  \draw[os arrow] (ir) -- (decode);
  \draw[os arrow] (decode) -- (registers);
  \draw[os arrow] (registers) |- (bus);
\end{pdfdiagram}
\diagramnote{箭头表示一次周期中的信息方向。CPU 先把 PC 放到地址侧；目标只在
ready 有效时交付数据。未映射窗口保持无响应，CPU 不能把全零误当成合法数据。}

这套指令编码刻意保持简单，所有 16-bit 立即数和地址都采用 little-endian：

\begin{osgridlongtable}{L{0.16\linewidth}L{0.25\linewidth}L{0.22\linewidth}L{0.27\linewidth}}
操作码 & 汇编形式 & 后续 byte & 状态变化 \\ \hline
\endfirsthead
操作码 & 汇编形式 & 后续 byte & 状态变化 \\ \hline
\endhead
\texttt{10} & \code{LOAD Rd, imm16} & 寄存器号、低 byte、高 byte &
  \(R_d\leftarrow imm\)，更新 Z \\ \hline
\texttt{20} & \code{ADD Rd, Rs} & 目标号、来源号 &
  \(R_d\leftarrow R_d+R_s\)，更新 Z \\ \hline
\texttt{30} & \code{STORE Rs, [addr16]} & 来源号、地址低 byte、地址高 byte &
  把寄存器低 byte 写入 RAM \\ \hline
\texttt{40} & \code{JMP addr16} & 地址低 byte、地址高 byte &
  \(PC\leftarrow addr\) \\ \hline
\texttt{41} & \code{JZ addr16} & 地址低 byte、地址高 byte &
  Z=1 时改写 PC，否则保留顺序 PC \\ \hline
\texttt{50} & \code{CALL addr16} & 地址低 byte、地址高 byte &
  压入返回 PC，再跳转 \\ \hline
\texttt{51} & \code{RET} & 无 &
  从栈弹出 PC \\ \hline
\texttt{FF} & \code{HALT} & 无 &
  停止继续取指 \\ \hline
\end{osgridlongtable}

\begin{workedexample}
ROM 从地址 \(\texttt{0x0000}\) 开始保存
\(\texttt{10 02 34 12}\)。复位后 \(PC=\texttt{0x0000}\)。

\begin{enumerate}[leftmargin=2.2em]
  \item 取地址 \(\texttt{0000}\) 的 \(\texttt{10}\)，PC 增为
    \(\texttt{0001}\)，译码为 \code{LOAD}；
  \item 取地址 \(\texttt{0001}\) 的 \(\texttt{02}\)，选择 R2；
  \item 取地址 \(\texttt{0002}\) 的 \(\texttt{34}\) 作为低 byte；
  \item 取地址 \(\texttt{0003}\) 的 \(\texttt{12}\) 作为高 byte；
  \item 合成为
    \(\texttt{0x34}\mathbin{|}(\texttt{0x12}\ll8)
      =\texttt{0x1234}\)；
  \item 写回后 R2 从 \(\texttt{0000}\) 变为
    \(\texttt{1234}\)，Z 保持 0，PC 已是 \(\texttt{0004}\)。
\end{enumerate}

这里有四次总线读取，却只有一次寄存器写回。把“执行一条指令”和“只过一个
周期”视为同一件事，会遗漏操作数读取和 ready 等待。
\end{workedexample}

\begin{center}
\small
\begin{osgridtabular}{L{0.08\linewidth}L{0.13\linewidth}L{0.11\linewidth}L{0.09\linewidth}L{0.09\linewidth}L{0.09\linewidth}L{0.10\linewidth}L{0.13\linewidth}}
周期 & 阶段 & 地址 & 数据 & ready & PC 后值 & R2 后值 & 发生的动作 \\ \hline
0 & fetch & 0000 & 10 & 1 & 0001 & 0000 & 取得操作码 \\ \hline
1 & decode & -- & -- & 1 & 0001 & 0000 & 选择 LOAD 状态 \\ \hline
2 & fetch & 0001 & 02 & 1 & 0002 & 0000 & 取得寄存器号 \\ \hline
3 & fetch & 0002 & 34 & 1 & 0003 & 0000 & 取得低 byte \\ \hline
4 & fetch & 0003 & 12 & 1 & 0004 & 0000 & 取得高 byte \\ \hline
5 & execute & -- & -- & 1 & 0004 & 0000 & 拼接立即数 \\ \hline
6 & writeback & -- & -- & 1 & 0004 & 1234 & 写 R2，更新 Z \\ \hline
\end{osgridtabular}
\end{center}

\topic{条件跳转让 PC 有两个候选值}

顺序执行时，取完 \code{JZ} 的三个 byte 后，PC 已经指向下一条指令。执行阶段
只需在“保留顺序 PC”和“写入跳转目标”之间选择。这个选择器就是前一章学过的
MUX，选择信号来自 Z。

\begin{pdfdiagram}
  \node[os state] (sequential) {顺序地址\\旧 PC + 指令长度};
  \node[os state,below=of sequential] (target) {目标地址\\指令中的 imm16};
  \node[os block,right=22mm of sequential,yshift=-10mm] (mux) {PC MUX\\选择信号 Z};
  \node[os state,right=22mm of mux] (next) {下一周期 PC};
  \draw[os arrow] (sequential) -- (mux);
  \draw[os arrow] (target) -- (mux);
  \draw[os arrow] (mux) -- (next);
\end{pdfdiagram}

示例程序先执行
\(\texttt{FFFF}+\texttt{0001}=\texttt{0000}\)。16-bit 结果回绕为零，
所以 Z 置 1。\code{JZ 0x0010} 取完后，顺序 PC 为
\(\texttt{0x000E}\)，目标为 \(\texttt{0x0010}\)：

\begin{osgridlongtable}{L{0.13\linewidth}L{0.16\linewidth}L{0.17\linewidth}L{0.17\linewidth}L{0.27\linewidth}}
时刻 & Z & 顺序 PC & 目标 & 写入 PC 的值 \\ \hline
\endfirsthead
时刻 & Z & 顺序 PC & 目标 & 写入 PC 的值 \\ \hline
\endhead
执行 ADD 前 & 0 & -- & -- & -- \\ \hline
ADD 写回后 & 1 & -- & -- & -- \\ \hline
JZ 取完操作数 & 1 & \texttt{000E} & \texttt{0010} & 尚未选择 \\ \hline
JZ 写回 & 1 & \texttt{000E} & \texttt{0010} & \texttt{0010} \\ \hline
\end{osgridlongtable}

若把 R1 的立即数改为 2，ADD 结果变成 \(\texttt{0001}\)，Z 为 0，JZ 不跳。
CPU 会在 \(\texttt{000E}\) 取到故意放置的 \code{HALT}。一处输入变化经过
ALU、标志寄存器和 PC MUX，最终改变了下一次总线地址。

\topic{CALL 和 RET 把返回地址保存到 RAM}

\code{CALL 0x0018} 占 3 byte，位于 \(\texttt{0x0010}\)。CPU 取完操作数
时，PC 已是 \(\texttt{0x0013}\)，这个值才是返回地址。初始
\(SP=\texttt{0xDFFE}\)，压入一个 16-bit 值的步骤为：

\[
SP\leftarrow SP-2=\texttt{0xDFFC}.
\]

x86 和本教学 CPU 都按 little-endian 保存多 byte 值，因此：

\[
M[\texttt{DFFC}]=\texttt{13},\qquad
M[\texttt{DFFD}]=\texttt{00}.
\]

\begin{center}
\begin{tikzpicture}[font=\footnotesize\sffamily]
  \foreach \y/\address/\before/\after in {
    3/DFFE/--/--,2/DFFD/00/00,1/DFFC/00/13,0/DFFB/00/00} {
    \node[anchor=east] at (0,\y) {\texttt{\address}};
    \draw[BookRule,fill=BookCodeBg] (0.3,\y-0.35) rectangle (2.0,\y+0.35);
    \node at (1.15,\y) {\texttt{\after}};
  }
  \draw[os arrow,BookRed] (3.7,2.8) -- (2.1,2.05);
  \node[anchor=west] at (3.8,2.8) {CALL 前 SP=\texttt{DFFE}};
  \draw[os arrow,BookTeal] (3.7,1.2) -- (2.1,1.05);
  \node[anchor=west] at (3.8,1.2) {CALL 后 SP=\texttt{DFFC}};
\end{tikzpicture}
\end{center}
\diagramnote{地址向上增大，栈向下增长。图中只画两个返回地址 byte；压栈不等于
把 SP 指向“最后一个空格”，SP 指向当前栈顶对象的低地址。}

\begin{workedexample}
子程序在 \(\texttt{0x0018}\) 装入 R2 后执行 \code{RET}。RET 先读
\(\texttt{DFFC}\) 得到低 byte \(\texttt{13}\)，再读
\(\texttt{DFFD}\) 得到高 byte \(\texttt{00}\)，合成
\(\texttt{0x0013}\)。随后
\[
PC\leftarrow\texttt{0x0013},\qquad
SP\leftarrow\texttt{0xDFFE}.
\]
地址 \(\texttt{DFFC}\) 和 \(\texttt{DFFD}\) 中的旧 byte 没有自动清零；
它们已经位于 SP 之下，因而不再属于当前栈内容。调试内存时要同时看 SP，不能
只凭残留 byte 判断栈里仍有一个返回地址。
\end{workedexample}

\topic{ready 延迟期间地址和控制必须保持稳定}

慢 ROM 可以在收到读请求后等待两个周期再返回数据。CPU 在等待期间不能递增
PC，也不能把地址换成下一个 byte：

\begin{center}
\begin{tikzpicture}[x=1.15cm,y=0.62cm,font=\footnotesize\sffamily]
  \foreach \y/\name in {5/CLK,4/address,3/read,2/ready,1/data} {
    \node[anchor=east] at (-0.3,\y) {\name};
    \draw[BookRule] (0,\y-0.35) -- (8,\y-0.35);
  }
  \foreach \x in {0,2,4,6} {
    \draw[BookBlue,thick] (\x,5) -- (\x,5.32) -- (\x+1,5.32) -- (\x+1,5);
  }
  \draw[BookTeal,thick] (0,4) -- (0.4,4.25) -- (6.0,4.25) -- (6.4,4);
  \node at (3.2,4.55) {\texttt{0x0018} 保持不变};
  \draw[BookBlue,thick] (0,3) -- (0.4,3.25) -- (6.0,3.25) -- (6.4,3);
  \draw[BookRed,thick] (0,2) -- (4.4,2) -- (4.8,2.25) -- (6.0,2.25) -- (6.4,2);
  \draw[BookAmber,thick] (0,1) -- (4.4,1) -- (4.8,1.25) -- (6.0,1.25) -- (6.4,1);
  \node at (5.4,1.55) {\texttt{0x10}};
  \draw[dashed,BookMuted] (4.8,0.4) -- (4.8,5.7);
  \node[below] at (4.8,0.4) {采样};
\end{tikzpicture}
\end{center}

模拟器把每次等待打印为 \code{ready=0}，最后一次响应打印为
\code{ready=1}。同一条正常程序在零等待与两拍等待下都应把
\(\texttt{0x34}\) 写入 \(\texttt{0x8000}\)，只有周期数增加。若等待期间
PC 提前递增，返回的 byte 会被错误地配给另一地址。

\lstinputlisting[
  language=C++,
  firstline=461,
  lastline=488,
  caption={只等待有限周期，超过上限就返回总线超时}
]{../../examples/teaching_cpu/source/teaching_cpu.cpp}

\topic{延迟、超时和没有目标是三种不同状态}

两拍 ready 延迟会在第三拍得到数据，属于正常的慢设备。模拟器还设置四拍
超时上限，并让另一个场景等待十二拍。CPU 只输出四行 \code{ready=0}，随后
结束为 \code{bus-timeout}；PC 保持在本次取指地址，没有接收尚未有效的数据。

\begin{osgridlongtable}{L{0.18\linewidth}L{0.23\linewidth}L{0.24\linewidth}L{0.25\linewidth}}
总线状态 & CPU 等待拍数 & 最终状态 & PC 与数据 \\ \hline
\endfirsthead
总线状态 & CPU 等待拍数 & 最终状态 & PC 与数据 \\ \hline
\endhead
ready 延迟 2 拍 & 2 & 继续运行 & 第 3 拍接收，随后 PC 递增 \\ \hline
ready 延迟 12 拍、上限 4 & 4 & \code{bus-timeout} & PC 不前进，不接收数据 \\ \hline
地址无目标 & 0 & \code{bus-no-response} & 地址译码失败，不等待 ready \\ \hline
\end{osgridlongtable}

这三个结果必须分开。把慢设备立即报成“无目标”，会错误拒绝合法事务；无限
等待慢设备，又会让整个宿主测试和单核系统永远停在同一条指令。

\topic{非法操作码和无响应目标必须结束当前路径}

ROM 的第一个 byte 若改成 \(\texttt{7E}\)，译码器找不到对应指令。模拟器把
状态改为 \code{illegal-opcode}，不猜测字节长度，也不继续取下一条。真实 CPU
同样需要为未定义指令进入明确异常路径；“忽略后继续”会让 PC 与指令边界失去
同步。

地址 \(\texttt{0xE000}\) 没有连接 ROM、RAM 或设备。向它执行 STORE 时，
模拟器打印一次 \code{ready=0} 并结束为 \code{bus-no-response}。硬件系统
可以选择总线错误、超时或机器检查，具体机制不同。软件至少要区分“目标返回
数值 0”“目标还没准备好”和“该地址没有目标”。

\begin{osgridlongtable}{L{0.21\linewidth}L{0.22\linewidth}L{0.23\linewidth}L{0.24\linewidth}}
故障 & 最后可信状态 & 不能继续的原因 & 可观察结果 \\ \hline
\endfirsthead
故障 & 最后可信状态 & 不能继续的原因 & 可观察结果 \\ \hline
\endhead
非法操作码 \texttt{7E} & PC 已越过该 byte & 不知道后续 byte 是否是操作数 &
  \code{illegal-opcode} \\ \hline
寄存器号 4 & 操作数已取完 & 机器只有 R0--R3 &
  \code{invalid-register} \\ \hline
读取 \texttt{E000} & 地址已驱动 & 没有目标交付数据 &
  \code{bus-no-response} \\ \hline
ready 超过四拍 & 地址和 read 保持稳定 & 等待预算已经用完 &
  \code{bus-timeout} \\ \hline
空栈执行 RET & SP 仍在初值 & 没有可读取的返回地址 &
  \code{stack-underflow} \\ \hline
无限跳转 & 状态仍可执行 & 教学程序不应永久占用宿主测试 &
  \code{instruction-limit} \\ \hline
\end{osgridlongtable}

\topic{运行 C++20 模拟器并对照正式实现}

可运行例子位于
\filepath{examples/teaching_cpu}。公共头文件定义 ISA 状态与总线边界，
源文件实现逐周期状态变化，\filepath{source/main.cpp} 组合五个运行场景。以下
片段就是正常程序的 ROM 字节；它不是只用于排版的伪代码：

\lstinputlisting[
  language=C++,
  firstline=67,
  lastline=124,
  caption={正常路径、非法操作码和未映射写的真实 ROM 字节}
]{../../examples/teaching_cpu/source/main.cpp}

在书稿根目录运行：

\begin{lstlisting}[language=bash]
python3 examples/teaching_cpu/verify.py
\end{lstlisting}

脚本使用 C++20、\code{-Wall -Wextra -Wpedantic -Werror} 编译，再检查正常、
ready 延迟、总线超时、非法操作码和总线无响应五种结果。构建发生在临时
目录，不会把宿主二进制混进书稿。

\begin{experimentbox}{跟踪 ADD、JZ、CALL 和 RET}
先运行验证脚本。把输出保存到文本文件后，按
\code{scenario=normal} 到下一条场景标记截取正常路径。

\begin{enumerate}[leftmargin=2.2em]
  \item 找到 ADD 的 writeback 周期，确认 R0 从 \texttt{FFFF} 变为
    \texttt{0000}，Z 同时从 0 变为 1；
  \item 找到 JZ 写回后的 PC，确认它从顺序候选
    \texttt{000E} 变为 \texttt{0010}；
  \item 找到两个 \code{phase=stack-write} 周期，记录地址和数据；
  \item 找到两个 \code{phase=stack-read} 周期，确认读回次序相同；
  \item 检查最后的 \code{memory[0x8000]=0x34}。
\end{enumerate}

每一步都同时记录周期号、PC、SP 和相关寄存器。这样才能说明“跳转发生”和
“返回地址确实经过 RAM”是两个不同结论。
\end{experimentbox}

\begin{faultinjectionbox}{交换返回地址的高低 byte}
在 \code{PushWord} 中交换 low byte 与 high byte 的写入位置。重新运行后，
RET 会把 \texttt{0x1300} 写入 PC，而不是 \texttt{0x0013}。下一次取指仍落在
ROM 范围内，但那里没有这段程序，常见结果是取到 \texttt{00} 并报告非法
操作码。

预期失败不能只写“程序不工作”。应记录：栈内两个 byte、RET 合成的目标 PC、
下一次取指地址和最终状态。恢复原顺序后再次运行，四种场景都应通过。
\end{faultinjectionbox}

\begin{pitfallbox}
\begin{itemize}[leftmargin=2.2em]
  \item PC 指向下一次取指地址；显示“当前指令地址”时要另存开始值；
  \item CALL 压入取完操作数后的 PC，不是 CALL 自身地址；
  \item ready=0 表示本周期不能接收数据，不表示数据值恰好为 0；
  \item 超时前每个等待周期都要保持地址、读写方向和 PC；
  \item 栈弹出后旧 byte 仍可留在 RAM，是否有效由 SP 决定；
  \item 非法操作码、未映射地址和正常 HALT 是三种不同停止原因。
\end{itemize}
\end{pitfallbox}

\topic{练习：自己增加一条指令}

\begin{enumerate}[leftmargin=2.2em]
  \item 手算 \(\texttt{0x7FFF}+\texttt{0x0001}\) 的 16-bit 结果和 Z。若新增
    signed overflow 标志，它应是什么？
  \item CALL 位于 \(\texttt{0x0200}\)，编码长度为 3，执行前
    \(SP=\texttt{0x9000}\)。写出执行后的 SP 和两个栈 byte 地址。
  \item ready 延迟从 0 改为 3。一次四 byte LOAD 至少增加多少个等待周期？
    只计算四次取 byte，不计译码和执行周期。
  \item 在当前 ISA 中加入 \code{LOADM Rd, [addr16]}。写出建议操作码、编码
    长度、读取顺序、失败状态和最小测试。
  \item 为什么不能把未映射地址的读取结果固定为 0，然后继续执行？
\end{enumerate}

\begin{answerbox}
\begin{enumerate}[leftmargin=2.2em]
  \item 结果为 \(\texttt{0x8000}\)，Z=0。两个正数相加得到符号位为 1 的
    结果，若按 16-bit signed 解释，overflow=1；
  \item 返回地址为 \(\texttt{0x0203}\)，新
    \(SP=\texttt{0x8FFE}\)。地址 \texttt{8FFE} 保存 \texttt{03}，
    \texttt{8FFF} 保存 \texttt{02}；
  \item 每个 byte 多等 3 拍，四次读取共增加 \(4\times3=12\) 拍；
  \item 可选 \texttt{11}，编码为操作码、寄存器号、地址低 byte、地址高
    byte。先验证寄存器号，再读取目标；目标无响应时保持寄存器旧值并报告
    \code{bus-no-response}。测试至少覆盖 RAM 成功读取和
    \texttt{E000} 失败；
  \item 0 是合法数据。混淆“返回 0”和“未响应”会让错误地址悄悄参与后续
    计算，故障位置被推迟，调试时只看到更远处的错误结果。
\end{enumerate}
\end{answerbox}
